%!TeX root=../emmatop.tex
\chapter[Chapter \thechapter]{}
\lettrine[lraise=0.3]{T}{he} appearance of the little sitting-room as they entered, was tranquillity itself; Mrs Bates, deprived of her usual employment, slumbering on one side of the fire, Frank Churchill, at a table near her, most deedily occupied about her spectacles, and Jane Fairfax, standing with her back to them, intent on her pianoforte.

\begin{figure}[tbph]
\centering
\includegraphics[width=.8\linewidth]{28spectacles}
\caption{Occupied about her spectacles}
\end{figure}

Busy as he was, however, the young man was yet able to shew a most happy countenance on seeing Emma again.

<This is a pleasure,> said he, in rather a low voice, <coming at least ten minutes earlier than I had calculated. You find me trying to be useful; tell me if you think I shall succeed.>

<What!> said Mrs Weston, <have not you finished it yet? you would not earn a very good livelihood as a working silversmith at this rate.>

<I have not been working uninterruptedly,> he replied, <I have been assisting Miss Fairfax in trying to make her instrument stand steadily, it was not quite firm; an unevenness in the floor, I believe. You see we have been wedging one leg with paper. This was very kind of you to be persuaded to come. I was almost afraid you would be hurrying home.>

He contrived that she should be seated by him; and was sufficiently employed in looking out the best baked apple for her, and trying to make her help or advise him in his work, till Jane Fairfax was quite ready to sit down to the pianoforte again. That she was not immediately ready, Emma did suspect to arise from the state of her nerves; she had not yet possessed the instrument long enough to touch it without emotion; she must reason herself into the power of performance; and Emma could not but pity such feelings, whatever their origin, and could not but resolve never to expose them to her neighbour again.

At last Jane began, and though the first bars were feebly given, the powers of the instrument were gradually done full justice to. Mrs Weston had been delighted before, and was delighted again; Emma joined her in all her praise; and the pianoforte, with every proper discrimination, was pronounced to be altogether of the highest promise.

<Whoever Colonel Campbell might employ,> said Frank Churchill, with a smile at Emma, <the person has not chosen ill. I heard a good deal of Colonel Campbell's taste at Weymouth; and the softness of the upper notes I am sure is exactly what he and all that party would particularly prize. I dare say, Miss Fairfax, that he either gave his friend very minute directions, or wrote to Broadwood himself. Do not you think so?>

Jane did not look round. She was not obliged to hear. Mrs Weston had been speaking to her at the same moment.

<It is not fair,> said Emma, in a whisper; <mine was a random guess. Do not distress her.>

He shook his head with a smile, and looked as if he had very little doubt and very little mercy. Soon afterwards he began again,

<How much your friends in Ireland must be enjoying your pleasure on this occasion, Miss Fairfax. I dare say they often think of you, and wonder which will be the day, the precise day of the instrument's coming to hand. Do you imagine Colonel Campbell knows the business to be going forward just at this time?—Do you imagine it to be the consequence of an immediate commission from him, or that he may have sent only a general direction, an order indefinite as to time, to depend upon contingencies and conveniences?>

He paused. She could not but hear; she could not avoid answering,

<Till I have a letter from Colonel Campbell,> said she, in a voice of forced calmness, <I can imagine nothing with any confidence. It must be all conjecture.>

<Conjecture—aye, sometimes one conjectures right, and sometimes one conjectures wrong. I wish I could conjecture how soon I shall make this rivet quite firm. What nonsense one talks, Miss Woodhouse, when hard at work, if one talks at all;—your real workmen, I suppose, hold their tongues; but we gentlemen labourers if we get hold of a word—Miss Fairfax said something about conjecturing. There, it is done. I have the pleasure, madam, (to Mrs Bates,) of restoring your spectacles, healed for the present.>

He was very warmly thanked both by mother and daughter; to escape a little from the latter, he went to the pianoforte, and begged Miss Fairfax, who was still sitting at it, to play something more.

<If you are very kind,> said he, <it will be one of the waltzes we danced last night;—let me live them over again. You did not enjoy them as I did; you appeared tired the whole time. I believe you were glad we danced no longer; but I would have given worlds—all the worlds one ever has to give—for another half-hour.>

She played.

<What felicity it is to hear a tune again which has made one happy!—If I mistake not that was danced at Weymouth.>

She looked up at him for a moment, coloured deeply, and played something else. He took some music from a chair near the pianoforte, and turning to Emma, said,

<Here is something quite new to me. Do you know it?—Cramer.—And here are a new set of Irish melodies. That, from such a quarter, one might expect. This was all sent with the instrument. Very thoughtful of Colonel Campbell, was not it?—He knew Miss Fairfax could have no music here. I honour that part of the attention particularly; it shews it to have been so thoroughly from the heart. Nothing hastily done; nothing incomplete. True affection only could have prompted it.>

Emma wished he would be less pointed, yet could not help being amused; and when on glancing her eye towards Jane Fairfax she caught the remains of a smile, when she saw that with all the deep blush of consciousness, there had been a smile of secret delight, she had less scruple in the amusement, and much less compunction with respect to her.—This amiable, upright, perfect Jane Fairfax was apparently cherishing very reprehensible feelings.

He brought all the music to her, and they looked it over together.—Emma took the opportunity of whispering,

<You speak too plain. She must understand you.>

<I hope she does. I would have her understand me. I am not in the least ashamed of my meaning.>

<But really, I am half ashamed, and wish I had never taken up the idea.>

<I am very glad you did, and that you communicated it to me. I have now a key to all her odd looks and ways. Leave shame to her. If she does wrong, she ought to feel it.>

<She is not entirely without it, I think.>

<I do not see much sign of it. She is playing Robin Adair at this moment—his favourite.>

Shortly afterwards Miss Bates, passing near the window, descried Mr Knightley on horse-back not far off.

<Mr Knightley I declare!—I must speak to him if possible, just to thank him. I will not open the window here; it would give you all cold; but I can go into my mother's room you know. I dare say he will come in when he knows who is here. Quite delightful to have you all meet so!—Our little room so honoured!>

She was in the adjoining chamber while she still spoke, and opening the casement there, immediately called Mr Knightley's attention, and every syllable of their conversation was as distinctly heard by the others, as if it had passed within the same apartment.

<How d' ye do?—how d'ye do?—Very well, I thank you. So obliged to you for the carriage last night. We were just in time; my mother just ready for us. Pray come in; do come in. You will find some friends here.>

So began Miss Bates; and Mr Knightley seemed determined to be heard in his turn, for most resolutely and commandingly did he say,

<How is your niece, Miss Bates?—I want to inquire after you all, but particularly your niece. How is Miss Fairfax?—I hope she caught no cold last night. How is she to-day? Tell me how Miss Fairfax is.>

And Miss Bates was obliged to give a direct answer before he would hear her in any thing else. The listeners were amused; and Mrs Weston gave Emma a look of particular meaning. But Emma still shook her head in steady scepticism.

<So obliged to you!—so very much obliged to you for the carriage,> resumed Miss Bates.

He cut her short with,

<I am going to Kingston. Can I do any thing for you?>

<Oh! dear, Kingston—are you?—Mrs Cole was saying the other day she wanted something from Kingston.>

<Mrs Cole has servants to send. Can I do any thing for you?>

<No, I thank you. But do come in. Who do you think is here?—Miss Woodhouse and Miss Smith; so kind as to call to hear the new pianoforte. Do put up your horse at the Crown, and come in.>

<Well,> said he, in a deliberating manner, <for five minutes, perhaps.>

<And here is Mrs Weston and Mr Frank Churchill too!—Quite delightful; so many friends!>

<No, not now, I thank you. I could not stay two minutes. I must get on to Kingston as fast as I can.>

<Oh! do come in. They will be so very happy to see you.>

<No, no; your room is full enough. I will call another day, and hear the pianoforte.>

<Well, I am so sorry!—Oh! Mr Knightley, what a delightful party last night; how extremely pleasant.—Did you ever see such dancing?—Was not it delightful?—Miss Woodhouse and Mr Frank Churchill; I never saw any thing equal to it.>

<Oh! very delightful indeed; I can say nothing less, for I suppose Miss Woodhouse and Mr Frank Churchill are hearing every thing that passes. And (raising his voice still more) I do not see why Miss Fairfax should not be mentioned too. I think Miss Fairfax dances very well; and Mrs Weston is the very best country-dance player, without exception, in England. Now, if your friends have any gratitude, they will say something pretty loud about you and me in return; but I cannot stay to hear it.>


\begin{figure}[tbph]
\centering
\includegraphics[width=.8\linewidth]{28moment}
\caption{<Oh! Mr Knightley, one moment more>}
\end{figure}

<Oh! Mr Knightley, one moment more; something of consequence—so shocked!—Jane and I are both so shocked about the apples!>

<What is the matter now?>

<To think of your sending us all your store apples. You said you had a great many, and now you have not one left. We really are so shocked! Mrs Hodges may well be angry. William Larkins mentioned it here. You should not have done it, indeed you should not. Ah! he is off. He never can bear to be thanked. But I thought he would have staid now, and it would have been a pity not to have mentioned.... Well, (returning to the room,) I have not been able to succeed. Mr Knightley cannot stop. He is going to Kingston. He asked me if he could do any thing....>

<Yes,> said Jane, <we heard his kind offers, we heard every thing.>

<Oh! yes, my dear, I dare say you might, because you know, the door was open, and the window was open, and Mr Knightley spoke loud. You must have heard every thing to be sure. <Can I do any thing for you at Kingston?> said he; so I just mentioned.... Oh! Miss Woodhouse, must you be going?—You seem but just come—so very obliging of you.>

Emma found it really time to be at home; the visit had already lasted long; and on examining watches, so much of the morning was perceived to be gone, that Mrs Weston and her companion taking leave also, could allow themselves only to walk with the two young ladies to Hartfield gates, before they set off for Randalls.